% Kapitel 3 -- Von Kuonrats vierzehntem Jahre und seiner Fahrt nach Hardegg

\section{Von Kuonr\^{a}ts vierzehntem Jahre}

\mylettrine{D}{er} genaue Tag, an welchem \gls{Kuonrat von Steinare} sein vierzehntes Jahr vollendete, ist mir nicht überliefert worden. Wohl aber sagen mehrere, die damals auf Burg Ottenstein dienten, dass die Wochen zuvor von solchem Regen erfüllt waren, dass die Wege im \gls{Waldviertel} zerflossen, die Hänge abrutschten und selbst die Jäger kaum weiter als bis an den nächsten Forstsaum gelangten. Also musste der Knabe auf der Burg verharren, der gern ritt, stritt und sich im Freien erprobte; und diese erzwungene Ruhe verdross ihn mehr, als Hunger oder Kälte es vermocht hätten.

Man meldet, er sei in jenen Tagen wie ein junger Hund im Hofe auf und ab gegangen, habe die Knechte mit unnützen Befehlen geplagt, die Hunde gereizt und an den nassen Mauern hinab in das Tal gestarrt, als könne sein Blick die Wolken selbst verjagen. Ob er dabei schon jenen zornigen Trotz in sich trug, der später sein Wesen so sehr zeichnete, vermag ich nicht mit Gewissheit zu sagen; doch passt die Überlieferung wohl zu dem Manne, der aus ihm ward.

Als endlich sein Tag gekommen war und der Regen eben soweit nachliess, dass man an Aufbruch denken durfte, liess \gls{Hadmar von Steinare} den Sohn vor sich treten. Dort wurde \gls{Kuonrat von Steinare} zum \gls{knappe}n bestimmt und dem Dienste seines Namensvetters, des Grafen \gls{Kuonrat von Hardegg}, zugewiesen, auf dass er bei einem mächtigeren Herrn lerne, was einem jungen Edelmann geziemt. Zugleich empfing er ein Ross, stark von Hals und Brust, unruhig im Auge und schnell im Tritt, welches fortan \gls{Veltloufer} hiess. Dass dieses Tier ihm lieber gewesen sei als mancher Christ, der ihm Gutes gönnte, will ich hier noch nicht behaupten; doch finden sich schon an diesem Anfang sichere Zeichen einer sonderbaren Zuneigung.

Als \gls{Kuonrat von Steinare} darauf gegen \gls{Hardegg} ritt, war das Land noch ganz vom Wasser gezeichnet. Die Räder hatten tiefe Furchen gezogen, in denen brauner Morast stund, und auch dort, wo kein Wagen gegangen war, sog der Weg an Huf und Schuh, als wolle er Mensch und Tier verschlingen. Neben ihm zog der \gls{Knecht}, und wer von beiden mehr unter der Mühsal litt, ist schwer zu sagen; denn das Ross ward von unten herauf bespritzt, der junge Herr bis an den Saum seines Wappenrockes beschmutzt, und der Diener trug denselben Schmutz auf Kleid, Händen und Gesicht.

Es begab sich aber, wie mir mehr als einmal erzählt ward, dass sie unterwegs an einer Schar junger Dorfmägde vorbeikamen, die am Rande des Weges standen und ihre Röcke vor dem Morast hoben. Da richtete \gls{Kuonrat von Steinare}, obgleich kaum erst dem Knabenalter entwachsen, sogleich Worte an sie, wie einer, der sich schon für gefällig bei Frauen hält. Er fragte eine nach ihrem Namen, einer andern rühmte er die hellen Haare, und einer dritten bot er, halb im Scherz, halb im Ernst, den Platz hinter sich auf dem Rosse an. Die Mägde aber sahen nicht auf das, was er sagen wollte, sondern auf das, was er war: ein schlammverkrusteter Knabe auf einem ebenso beschmierten Tiere. Da lachten sie, zuerst leise, dann offen, und eine soll gerufen haben: \enquote{Herr, lasst erst Euren Gaul den Fluss finden, ehe Ihr nach Mädchen greift.}

Solcher Hohn traf ihn hart. Nicht gegen die Mägde wandte sich sein Zorn, denn ihr Lachen mochte er noch eher als Trotz deuten, sondern gegen den \gls{Knecht}, der ihm wie immer am nächsten stand. Er schalt ihn faul und unachtsam und gab ihm die Schuld daran, dass Ross und Rock in solchem Zustande seien, als hätte ein armer Diener den Regen selbst befehlen können. Von da an, so heisst es, befahl er, der \gls{Knecht} solle alle paar Stunden \gls{Veltloufer} und den Wappenrock mit Wasser, Stroh und Tuch reinigen, sooft sie rasteten oder an einen Bach gelangten. Darin zeigt sich schon früh jene Art von Hochmut, die nicht fragt, was möglich sei, sondern nur, was dem eigenen Dünkel frommt.